% =============================================================================
%  
%  
% =============================================================================

\section{Avances (Antes - Después)}

% ── Tabla comparativa en tres diapositivas ────────────────────────────────────
\begin{frame}{Avances de 2026-2 a 2027-1}
  \begin{table}
    \centering
    \footnotesize
    \setlength{\tabcolsep}{3pt}
    \begin{tabularx}{\linewidth}{@{} >{\raggedright\arraybackslash}p{0.18\linewidth} *{3}{>{\raggedright\arraybackslash}X} @{}}
      \toprule
      \textbf{Eje de análisis}   & \textbf{Sem 2026-2} & \textbf{Sem 2027-1} & \textbf{Estado final} \\
      \midrule
      Stack de Software             & \termino{Apimetro} en producción,
                                      \termino{VFTModel} en desarrollo de fases 2 y 3 sin terminar.
                                    & Capas de Software operables.
                                    & Terminado para la tesis, liberado para la comunidad. \\

      Grafo Multimodal              & En proceso de integración sobre rutas principales
                                    & Grafo consolidado con \textbf{11\,115} nodos y
                                    \textbf{45\,517} aristas
                                    con 12 sistemas de transporte
                                    (MB, Metro, Trolebús, Cbb, CC, RTP, Pumabús, etc.)
                                    & Suficiente para la tesis, pero valdría la pena revisar. \\
      \bottomrule
    \end{tabularx}
  \end{table}
\end{frame}

\begin{frame}{Avances de 2026-2 a 2027-1}
  \begin{table}
    \centering
    \footnotesize
    \setlength{\tabcolsep}{3pt}
    \begin{tabularx}{\linewidth}{@{} >{\raggedright\arraybackslash}p{0.18\linewidth} *{3}{>{\raggedright\arraybackslash}X} @{}}
      \toprule
      \textbf{Eje de análisis}   & \textbf{Sem 2026-2} & \textbf{Sem 2027-1} & \textbf{Estado final} \\
      \midrule
      Impedancia y tiempo           & Suposición teórica inicial con 85~m de tolerancia de traslado.
                                    & Algoritmo de \termino{Dijkstra} con fricción vial y
                                    penalización por transbordo ($T_b$). En la red base, $T = 113.91$~min.
                                    & Suficiente para la tesis. Valdría la pena ejecutar
                                    otros algoritmos de traslado para otro estudio. \\

      Simulación del Anillo         & Propuesta exploratoria.
                                    & Simulación en los cuatro cuadrantes.
                                    Caso~1, Periférico vial.
                                    Caso~2, Periférico con Metrobús.
                                    Caso~3, Periférico con Metro.
                                    & Se crearon dos casos independientes, que amplían el alcance de la tesis.
                                    Para mejores estudios se pueden simular más casos. \\
      \bottomrule
    \end{tabularx}
  \end{table}
\end{frame}

\begin{frame}{Avances de 2026-2 a 2027-1}
  \begin{table}
    \centering
    \footnotesize
    \setlength{\tabcolsep}{3pt}
    \begin{tabularx}{\linewidth}{@{} >{\raggedright\arraybackslash}p{0.18\linewidth} *{3}{>{\raggedright\arraybackslash}X} @{}}
      \toprule
      \textbf{Eje de análisis}   & \textbf{Sem 2026-2} & \textbf{Sem 2027-1} & \textbf{Estado final} \\
      \midrule
      Síntesis                      & Formulación conceptual de indicadores aislados.
                                    & Se crea el Espectro MJG para simplificar el diagnóstico
                                    por bandas de escala; dando como resultado la clasificación Garibelt.
                                    & Excedente para la tesis pero necesario. Valdría la pena demostrar
                                    teóricamente y matemáticamente si la clasificación puede pasar
                                    a un indicador universal. \\
      \bottomrule
    \end{tabularx}
  \end{table}
\end{frame}
